\section{FR Representations and Equivariant Transformations}
\label{sec:operators}
This section does not study screening directly. Instead, it introduces FR representations together with the primitive transformations acting on them: FR duality, translation, and restriction. We also establish that translation and restriction are equivariant under FR duality.

Throughout this section, functions take values in the extended reals \(\overline{\mathbb R}:=\mathbb R\cup\{+\infty\}\), so that a domain constraint can be encoded as \(+\infty\) outside the feasible set. For a finite set \(N\) and \(I\subseteq N\), \(\bar I:=N\setminus I\) denotes its complement, and  disjoint-union \(I\sqcup\bar I=N\) records that \(I\) and \(\bar I\) partition \(N\). This disjoint-union notation is used for both index sets fixed below: \(I\sqcup\bar I=\{1,\ldots,n\}\) and \(J\sqcup\bar J=\{1,\ldots,m\}\).






%%%%%%%%%%%%%%%%%%%%%%%%%%%%%%%%%%%%%%%%%%%%%%%%%%%%%%%%%%%%%%%

\subsection{FR Representations}

We consider the problems of the form
\begin{equation}
\min_{x \in \mathbb R^n}\;
p(x), \qquad p(x) = f(Ax)+g(x)+a.
\label{eq:FR-problem}
\end{equation}
where
\begin{equation}
f:\mathbb R^m\rightarrow\overline{\mathbb R},
\qquad
g:\mathbb R^n\rightarrow\overline{\mathbb R},
\qquad
A\in\mathbb R^{m\times n},
\qquad
a\in\mathbb R,
\label{eq:cond-of-FR-representation}
\end{equation}
with \(f\) and \(g\) closed, proper, and convex.
We further assume that \(f\) and \(g\) are separable, \(f(y)=\sum_{j=1}^m f_j(y_j)\)
and \(g(x)=\sum_{i=1}^n g_i(x_i)\), with each \(f_j\) and \(g_i\) closed, proper, and convex on $\mathbb{R}$. Here, at the first glance, \(a\) is a redundant constant since it does not affect the minimizer of \(p\). However, it is important to keep track of \(a\) because it is used as an absorbing constant in the following transformations. 


\noteLE{add ref}.
The problem \eqref{eq:FR-problem} is of particular importance in both
convex optimization and machine learning. In convex optimization, it
forms the foundation of FR duality theory. In machine learning, it
provides a unified formulation for a wide range of learning problems,
including regression, classification, and optimal transport.


\begin{definition}[FR representation]
    If \((f,g,A,a)\) satisfies \eqref{eq:cond-of-FR-representation}, then we say that it is well-defined. In this case, we refer to it as an FR representation of function \(p\).
\end{definition}

In this paper, we identify the optimization problem \eqref{eq:FR-problem} with its objective function \(p\) and its FR representation \((f,g,A,a)\).
An FR representation is not merely a convenient way to write down a function. The quadruple \((f,g,A,a)\) is designed so that the class of FR representations is \emph{closed} under the three transformations studied in this section. 

For function \(p\) in \eqref{eq:FR-problem}, we write \(\operatorname{size}(p):=(m,n)\) the \emph{problem size}. We also refer to it as the \emph{intrinsic complexity} of the problem. 
We define $\Gamma_{m, n}$ to be the set of all functions $p$ with problem size $(m, n)$, and $\Gamma$ to be the set of all FR representations of any size.
If both \(f\) and \(g\) are separable, then \((f,g,A,a)\) is said to be a \emph{separable FR representation}. Intrinsic complexity ans separability are essential and will be discussed in more details in the next section.

%%%%%%%%%%%%%%%%%%%%%%%%%%%%%%%%%%%%%%%%%%%%%%%%%%%%%%%%%%%%%%%





\subsection{FR Duality}

% main message: the dual of an FR representation is another FR representation, and duality is an involution on the class of FR representations.

For a function \(\varphi:\mathbb R^d\to\overline{\mathbb R}\), its Fenchel
conjugate is
\[
\varphi^*(y):=\sup_{x\in\mathbb R^d}\{\langle y,x\rangle-\varphi(x)\}.
\]
When \(\varphi\) is closed, proper, and convex, then
\(\varphi^{**}=\varphi\), a fact used repeatedly below. \noteLE{add ref.}

The FR dual problem of \eqref{eq:FR-problem} is
\[
\max_{y \in \mathbb R^m}\;
-d(y), \qquad d(y) = g^*(-A^\top y) + f^*(y)-a.
\]
Note that in the classical FR duality, we do not have \(a\) and \(-a\) in \(p\) and \(d\). 
Here the constant \(a\) is negated in the dual, so that the weak duality inequality \(p(x) + d(y)\geq 0\) remains valid for all \(x\in \mathbb R^n\) and \(y \in \mathbb R^m\).

The dual function $d$ is itself an FR representation, with the problem size \((n, m)\). Therefore, FR duality can also be considered as a transformation on the class of FR representations. This motivates the following definition.


\begin{definition}[FR Duality]
The FR duality is a transformation \((\cdot)^*: \Gamma_{m, n} \to \Gamma_{n, m}\) such that the image of \(p=(f,g,A,a)\) is
\[
p^*
:=
(g^*,\,f^*,\,-A^\top,\,-a).
\]
\end{definition}

Note that FR duality is also a transformation on the class of FR representations $\Gamma$.




\noteCLAUDE{This result is used for simultaneous screening.} \noteLE{Verified.}
\begin{proposition}[FR Duality Involution]
\label{prop:involution}
For every FR representation \(p \in \Gamma\), we have \((p^*)^*=p\).
\end{proposition}

\begin{proof}
Write \(p=(f,g,A,a)\), so \(p^*=(g^*,f^*,-A^\top,-a)\). Applying the
same rule again,
\[
(p^*)^*
=
\bigl((f^*)^*,\,(g^*)^*,\,-(-A^\top)^\top,\,-(-a)\bigr)
=
(f^{**},g^{**},A,a)
=
(f,g,A,a)
=
p,
\]
using biconjugation \(f^{**}=f\), \(g^{**}=g\), which hold true since $f$ and $g$ are closed, proper, and convex.
\end{proof}


FR duality is itself a transformation on \(\Gamma\). 
It naturally induces a notion of duality between transformations, called \emph{equivariance}, this is the main structure to investigate in the remaining paper.

\begin{definition}[Equivariance under FR duality]
\label{def:equivariance}
Transformations \(F,G:\Gamma\to\Gamma\) are \emph{equivariant under FR duality} if
\[
(F(p))^*=G(p^*),
\qquad
\forall\,p\in\Gamma.
\]
\end{definition}









%%%%%%%%%%%%%%%%%%%%%%%%%%%%%%%%%%%%%%%%%%%%%%%%%%%%%%%%%%%%%%%








\subsection{Translation}
\label{subsec:translation}
We first introduce three premitive translations.

\begin{definition}[Primitive translations]
For \(b\in\mathbb R^d\), \(c\in\mathbb R\), and \(d\in\mathbb R^d\),
define, for a function \(\varphi: \mathbb R^n \rightarrow \overline{\mathbb R}\),
\[
\text{domain translation:}
\quad
(T^{\mathrm d}_b \varphi)(y):=\varphi(y+b),
\]
\[
\text{value translation:}
\quad
(T^{\mathrm v}_c \varphi)(y):=\varphi(y)+c,
\]
\[
\text{slope translation:}
\quad
(T^{\mathrm s}_d\varphi)(x):=\varphi(x)+\langle d,x\rangle.
\]
\end{definition}

Geometrically, domain translation shifts the graph horizontally,
value translation shifts it vertically, and slope translation tilts
the graph by adding a linear function, thereby shifting every
subgradient by the fixed vector \(d\).


\begin{definition}[Translation parameter and its space]
\label{def:translation-parameter}
We call
\(
K_{m,n}:=\mathbb R^m\times\mathbb R\times\mathbb R^n,
\)
a \emph{translation parameter space}.
An element \(k=(b,c,d)\in K_{m,n}\) is called a \emph{translation
parameter}.
\end{definition}

Now, we can define a translation indexed by a translation parameter.

\begin{definition}[Translation]
Let \(p=(f,g,A,a)\in \Gamma_{m, n}\). The translation $T_k$ indexed by \(k=(b,c,d) \in K_{m, n}\) on \(p\) is defined by \(T_k : \Gamma_{m, n} \rightarrow \Gamma_{m, n} \) such that 
\[
T_k(p)
:=
(T^{\mathrm d}_bf,\;
T^{\mathrm s}_dg,\;
A,\;
a+c).
\]
\end{definition}


Note that, the matrix \(A\) is untouched by translation: only the two functions
and the constant move. This is what lets restriction and translation
later be composed without the matrix itself being disturbed by which
primitive acts first.

It is clear that the family \(\{T_k\}\) forms an Abelian translation action on the class of FR representations:
\[
T_{k_1+k_2}(p)
=
T_{k_1}(T_{k_2}(p)) = T_{k_2}(T_{k_1}(p)). 
\]
\noteCLAUDE{Not currently cited by any result} \noteLE{verified. But kept because of its beauty. So keept it of this form, do not replace it by a proposition. The commutative will be used in the next paper.}

Section~\ref{sec:screening} will apply the translation action after
choosing \(k\) as a function of the coordinates being screened. Before
that, we establish this subsection's law: how the translation action
interacts with FR duality.



\begin{proposition}[Primitive Translation Equivariance]
\label{prop:duality-of-primitive-translations}
For every convex function $\varphi: \mathbb{R}^n \to \overline{\mathbb{R}}$, $b,d \in \mathbb{R}^n$, and $c \in \mathbb{R}$, we have 
\begin{enumerate}
    \item \((T^{\mathrm d}_b \varphi)^* = T^{\mathrm s}_{-b}(\varphi^*)\)
    \item \((T^{\mathrm s}_d \varphi)^* = T^{\mathrm d}_{-d}(\varphi^*)\)
    \item \((T^{\mathrm v}_c \varphi)^* = T^{\mathrm v}_{-c}(\varphi^*)\)
\end{enumerate}
\end{proposition}
\begin{proof}
    \noteLE{TODO. Add ref.}
\end{proof}
\noteCLAUDE{Used in
in the proof of Translation Equivariance below.}\noteLE{Verified.}


We define the \emph{dual parameter} of \(k\) by
\(k^*=(-d,-c,-b)\in K_{n, m}\).
Intuitively, the duality map \(k \mapsto k^*\) swaps the domain and slope translation parameters, and negates all three parameters.
Thus, $*: K_{m, n} \to K_{n, m}$ is a duality map such that \((k^*)^*=k\).

\begin{proposition}[Translation Equivariance]
\label{thm:fenchel-equivariance}
Let \(p\in \Gamma_{m, n}\) be an FR representation and let \(k=(b,c,d)\in K_{m, n}\) be a translation
parameter. Then
\[
(T_k(p))^*
=
T_{k^*}(p^*).
\]
\end{proposition}

\paragraph{Proof.}
Write \(p=(f,g,A,a)\), so \(T_k(p)=(T^{\mathrm d}_bf,\,T^{\mathrm
s}_dg,\,A,\,a+c)\). Applying the FR dual, which negates the constant
slot, \((f_0,g_0,A_0,a_0)^*=(g_0^*,f_0^*,-A_0^\top,-a_0)\), termwise, and
using Proposition~\ref{prop:duality-of-primitive-translations},
\[
(T_k(p))^*
=
\big(
(T^{\mathrm s}_dg)^*,\;
(T^{\mathrm d}_bf)^*,\;
-A^\top,\;
-(a+c)
\big)
=
\big(
T^{\mathrm d}_{-d}(g^*),\;
T^{\mathrm s}_{-b}(f^*),\;
-A^\top,\;
-a-c
\big).
\]

On the other side, \(p^*=(g^*,f^*,-A^\top,-a)\), and \(k^*=(-d,-c,-b)\),
so \(T_{k^*}(p^*)\) applies a domain translation by \(-d\) to \(p^*\)'s
first slot, a slope translation by \(-b\) to its second slot, and adds
\(-c\) to its constant:
\[
T_{k^*}(p^*)
=
\big(
T^{\mathrm d}_{-d}(g^*),\;
T^{\mathrm s}_{-b}(f^*),\;
-A^\top,\;
-a-c
\big).
\]
Every term matches \((T_k(p))^*\), which proves the theorem.
\qed


\noteCLAUDE{Used in  Section~\ref{sec:screening}'s Feature--Sample
Duality theorem.}
\noteLE{Verified.}
%%%%%%%%%%%%%%%%%%%%%%%%%%%%%%%%%%%%%%%%%%%%%%%%%%%%%%%%%%%%%%%









%%%%%%%%%%%%%%%%%%%%%%%%%%%%%%%%%%%%%%%%%%%%%%%%%%%%%%%%%%%%%%%








\subsection{Restriction}


Let \(I\subseteq\{1,\ldots,n\}\) and \(J\subseteq\{1,\ldots,m\}\). For a
separable function \(g(x)=\sum_{i=1}^n g_i(x_i)\), write
\(x_I:=(x_i)_{i\in I}\) and \(g_I(x_I):=\sum_{i\in I}g_i(x_i)\); define
\(f_J\) symmetrically for a separable \(f(y)=\sum_{j=1}^m f_j(y_j)\).
For a matrix \(A\in\mathbb R^{m\times n}\), write \(A_I:=A_{:,I}\) for
the submatrix of columns indexed by \(I\), and \(A_J:=A_{J,:}\) for the
submatrix of rows indexed by \(J\). Combining both,
\(A_{J,I}:=(A_J)_I=(A_I)_J\) denotes the submatrix with rows in \(J\)
and columns in \(I\).


\begin{definition}[Feature and sample restriction]
Let \(I\sqcup\bar I=\{1,\ldots,n\}\) and \(J\sqcup\bar J=\{1,\ldots,m\}\). Define the feature restriction operator $R^{F}_I:\Gamma_{m,n}\to\Gamma_{m,|I|}$ and the sample restriction operator \(R^{S}_J:\Gamma_{m,n}\to\Gamma_{|J|,n}\) by
\[
R^{F}_I(f,g,A,a)
:=
(f,g_I,A_I,a).
\]
and 
\[
R^{S}_J(f,g,A,a)
:=
(f_J,g,A_J,a).
\]
\end{definition}

Feature restriction only touches \(g\) and the columns of \(A\); sample
restriction only touches \(f\) and the rows of \(A\). Because they act
on disjoint parts of \(p\), the order in which they are applied cannot
matter: this subsection's law.


\begin{proposition}[Restriction Equivariance]
\label{prop:restriction-equivariance}
For every \(p=(f,g,A,a)\in\Gamma\), every \(I\subseteq\{1,\ldots,n\}\),
and every \(J\subseteq\{1,\ldots,m\}\),
\[
\bigl(R^{F}_I(p)\bigr)^*=R^{S}_I(p^*),
\qquad
\bigl(R^{S}_J(p)\bigr)^*=R^{F}_J(p^*).
\]
\end{proposition}

\begin{proof}
Write \(p^*=(g^*,f^*,-A^\top,-a)\). Applying the FR dual termwise to
\(R^F_I(p)=(f,g_I,A_I,a)\),
\[
\bigl(R^F_I(p)\bigr)^*
=
\bigl((g_I)^*,\;f^*,\;-A_I^\top,\;-a\bigr).
\]
By separability of \(g\), \((g_I)^*=(g^*)_I\). Rows \(I\) of
\(-A^\top\) equal \(-A_I^\top\), since rows of \(A^\top\) are columns of
\(A\) transposed. Hence
\[
\bigl(R^F_I(p)\bigr)^*
=
\bigl((g^*)_I,\;f^*,\;-A_I^\top,\;-a\bigr)
=
R^S_I(p^*),
\]
the last equality by the definition of sample restriction applied to
\(p^*=(g^*,f^*,-A^\top,-a)\). The second identity is the mirror
computation, restricting \(R^S_J(p)=(f_J,g,A_J,a)\) instead, and using
\((f_J)^*=(f^*)_J\) by the same separability argument.
\end{proof}

\noteCLAUDE{Restriction Equivariance is used in the proof of the Feature--Sample Duality theorem in Section~\ref{sec:screening}.} \noteLE{Verified.}
